\section{\code{Money}, and the Scalar That Promises Nothing}
\label{sec:ch13-money-and-the-scalar}
\index{value type!declared in two subgraphs|(}%

\code{Money} is two fields and is declared twice, in Pricing and in Ordering,
\code{@shareable} in both. It is not an entity, has no key and cannot be fetched:
it is a shape that two services independently know how to fill in. Federation
calls that a value type, and Mosaic has had one since
chapter~\ref{ch:strangling-the-monolith} without anybody testing what holds it
together.

Three edits, three different outcomes, and they get progressively quieter.

\subsection*{An extra field: caught, and not for the reason you would guess}

Give Ordering's copy a field Pricing's does not have:

\begin{graphqlsdl}[fontsize=\footnotesize]
type Money @shareable { amount: Decimal!  currency: String!  formatted: String! }
\end{graphqlsdl}

\begin{minted}[fontsize=\scriptsize,breaklines]{text}
The field "formatted" is unresolvable at the following path:
 query {
  products {
   price {
    formatted <--
   }
  }
 }
This is because:
 - The root type field "Query.products" is defined in the following subgraph: "catalog".
 - The field "Money.formatted" is defined in the following subgraph: "ordering".
 - The entity ancestor "Product" in subgraph "pricing" has no accessible target entities (resolvable @key directives)
in the subgraphs where "Money.formatted" is defined.
 - The type "Money" is not a descendant of any other entity ancestors that can provide a shared route to access
"formatted".
 - The type "Money" has no accessible target entities (resolvable @key directives) in any other subgraph, so accessing
other subgraphs is not possible.
\end{minted}

\index{satisfiability!and value types}%
Not one word of that is about value types. It is chapter~\ref{ch:composition}'s
satisfiability walk, arriving at a field it cannot reach and listing every route
it tried. The last two lines are the finding: \code{Money} has no key, so there
is no representation the router could send to go and get \code{formatted}. A
value type has to be identical in every subgraph that declares it, and the reason
is not a rule about value types. It is that the graph has no way to fetch part of
one.

\subsection*{A missing \code{@shareable}: caught, and specific}

Take the directive off Ordering's copy:

\begin{minted}[fontsize=\scriptsize,breaklines]{text}
The Object "Money" defines the same fields in multiple subgraphs without the "@shareable" directive:
 The field "amount" is defined and declared "@shareable" in the following subgraph: "pricing".
 However, it is not declared "@shareable" in the following subgraph: "ordering".
 The field "currency" is defined and declared "@shareable" in the following subgraph: "pricing".
 However, it is not declared "@shareable" in the following subgraph: "ordering".
\end{minted}

Chapters~\ref{ch:composition} and~\ref{ch:strangling-the-monolith} between them
have four different mistakes that report as a shareability error, and this is the
one case where the error is about what it says it is about. Worth recording,
because it means the message is not useless: it is over-used.

\subsection*{A nullability drift: not caught at all}

Now change one character. Ordering declares \code{currency: String} where Pricing
declares \code{currency: String!}:

\begin{graphqlsdl}[fontsize=\footnotesize]
type Money { amount: Decimal!  currency: String }
\end{graphqlsdl}

That is the composed client schema. It composes, and the weaker declaration wins
for every client of the graph.

\index{nullability!the more permissive side wins}%
Chapter~\ref{ch:entity-resolution-done-right} found this exact behaviour for an
\code{@external} field and recorded it as a decision: an \code{@external} copy is
a copy of somebody else's contract and has to match it. Here it is again for a
shared value type, and the second sighting is what makes it a rule rather than an
anecdote. \textbf{Where two subgraphs declare the same thing and differ only in
strictness, composition takes the looser side and says nothing.} A team that
relaxes a field in their own service relaxes it in everybody's.

\subsection*{And a scalar, where there is nothing to check}

\index{scalar!has no structure to check}%
\code{Decimal} is declared in Pricing and in Ordering, both with a
\code{@specifiedBy} url pointing at ChilliCream's definition. Point one of them
somewhere else entirely:

\begin{graphqlsdl}[fontsize=\footnotesize]
scalar Decimal @specifiedBy(url: "https://example.com/mosaic/cents-as-an-integer")
\end{graphqlsdl}

The graph composes. There is no message. And the composed client schema is this:

\begin{graphqlsdl}[fontsize=\footnotesize]
scalar Decimal
\end{graphqlsdl}

Both urls are dropped on the way in, so a client introspecting the graph cannot
even see that there was a disagreement to have. Two teams can mean euros and
cents by the same name and every tool in this stack will agree that the graph is
fine.

\index{composition!what it can and cannot check}%
Which is the whole point of putting these three edits next to each other. The
composer is a structural tool and its protection tracks how much structure a type
has, exactly. An object type has fields, so a missing field is unreachable and a
weakened field is merged permissively and silently. An enum has members, so its
members are checked against the direction they travel. A scalar has a name and
nothing else, so a scalar is checked for having a name.

That is not a defect to be fixed in a later version. Nothing could check it: the
meaning of a custom scalar exists in two codebases and in no document either tool
can read. What follows is a working rule rather than a complaint. Every shared
type in a federated graph needs an owner and a written definition somewhere
outside both schemas, and the amount of that written definition you need is
inversely proportional to how much structure the type has. \code{Money} needs a
sentence. \code{Decimal} needs a paragraph about precision and rounding, and the
teams have to read it, because nothing else in this toolchain ever will.

\medskip
Four problems, and only one of them was solved rather than described. That is a
fair summary of the chapter and of the subject: a federated graph makes the
shapes agree and cannot make the meanings agree, and the second half is a
question for people. Chapter~\ref{ch:real-time-in-a-federated-world} takes the
graph this chapter leaves and asks what happens when the answer arrives more than
once.

\index{value type!declared in two subgraphs|)}%
